
\documentclass[a4paper,12pt]{article}
\usepackage[italian]{babel}
\begin{document}

\title{Verifica di Basi di Dati – SQL}
\date{27/03/2026}
\maketitle

Dato lo schema logico di seguito descritto, scrivi il codice SQL richiesto. Non è consentito usare EXIST, ANY, LIMIT.

\section*{Schema Logico}
CLIENTI(id cliente, nome, cognome, città, patente)\\
AUTO(id auto, marca, modello, targa, anno, id categoria)\\
CATEGORIE(id categoria, nome categoria, costo giornaliero)\\
NOLEGGI(id noleggio, id cliente, id auto, data inizio, data fine)\\
PAGAMENTI(id pagamento, id noleggio, importo, data pagamento)\\
MANUTENZIONI(id manutenzione, id auto, data manutenzione, costo)

\section*{Parte 1 – Query SQL}
Scrivere le query SQL per le seguenti richieste.
\begin{enumerate}
\item Elencare nome, cognome e città dei clienti residenti a Roma o Milano, ordinati per cognome.
\item Elencare marca, modello e costo giornaliero delle auto appartenenti alla categoria “SUV”.
\item Elencare i noleggi con nome cliente, targa auto, data inizio e data fine.
\item Elencare il numero di auto per ogni categoria.
\item Elencare i clienti che hanno effettuato almeno un noleggio.
\item Elencare i clienti che hanno effettuato più di 2 noleggi.
\item Calcolare l’importo totale pagato per ogni noleggio.
\item Calcolare la spesa totale per ogni cliente.
\item Elencare le auto che hanno avuto almeno una manutenzione.
\item Elencare le auto che non sono mai state noleggiate.
\item Elencare il numero di noleggi per ogni auto.
\item Elencare i clienti che hanno noleggiato auto di categoria “Lusso”.
\item Elencare le categorie con costo giornaliero medio superiore a 50.
\item Trovare l’auto più noleggiata.
\end{enumerate}

\section*{Parte 2 – INSERT}
Inserire un nuovo cliente: Luca Bianchi, residente a Torino, patente AB123CD.

\section*{Parte 3 – UPDATE}
Aggiornare il costo giornaliero della categoria “SUV” aumentandolo del 10\%.

\section*{Parte 4 – SQL DDL}
Scrivere le istruzioni SQL per creare le seguenti tabelle:

UFFICI(id ufficio, città, indirizzo)\\
DIPENDENTI(id dipendente, nome, cognome, id ufficio)

Vincoli:
\begin{itemize}
\item id\_ufficio e id\_dipendente chiavi primarie
\item città non nulla
\item id\_ufficio in DIPENDENTI chiave esterna
\item nome e cognome non nulli
\end{itemize}

\section*{Parte 5 – Viste}
Creare una vista chiamata spesa\_clienti che mostri:
\begin{itemize}
\item id\_cliente
\item nome
\item cognome
\item totale speso
\end{itemize}

\end{document}
